%% Part 3: Comparative Analysis and Contributions
%% NeurIPS 2026 positioning paper
%% Depends on: Part 1 (\ref{sec:related}), Part 2 (method/formulation sections)

\section{Comparative Analysis and Contributions}
\label{sec:comparison}

With the related work surveyed in Section~\ref{sec:related}, we now
position our work against prior art.
We first provide a structured comparison across key desiderata
(Section~\ref{sec:comparison_table}), then present focused head-to-head analyses
with each method category (Section~\ref{sec:head_to_head}), situate our
empirical results against reference baselines (Section~\ref{sec:results_position}),
enumerate our concrete contributions (Section~\ref{sec:contributions}),
and conclude with an honest assessment of the current scope and limitations
(Section~\ref{sec:limitations}).


%% ============================================================
\subsection{Systematic Comparison}
\label{sec:comparison_table}

Table~\ref{tab:comparison} compares representative methods across eight
dimensions that collectively define the desiderata for clinical time-series
domain adaptation with a frozen predictor.
No prior method satisfies all eight.

\begin{table}[t]
\centering
\caption{Systematic comparison of domain adaptation and frozen-model methods
across key desiderata for clinical time-series translation.
\textbf{Model Frozen}: whether the target predictor's weights are entirely
fixed (\cmark), partially fixed (P), or trainable (\xmark).
\textbf{Adapt. Space}: Input (I), Representation (R), Weight (W), or Output (O).
\textbf{Sample-Cond.}: whether the transformation depends on the input instance.
\textbf{Clinical Tasks}: number of distinct clinical prediction tasks evaluated.
}
\label{tab:comparison}
\small
\setlength{\tabcolsep}{3.2pt}
\begin{tabular}{@{}lcccccccc@{}}
\toprule
\textbf{Method} &
\rotatebox{60}{\textbf{Model Frozen?}} &
\rotatebox{60}{\textbf{Adapt. Space}} &
\rotatebox{60}{\textbf{Sample-Cond.?}} &
\rotatebox{60}{\textbf{Retrieval?}} &
\rotatebox{60}{\textbf{Var.-Len. TS?}} &
\rotatebox{60}{\textbf{Clinical Tasks}} &
\rotatebox{60}{\textbf{Grad. Theory?}} \\
\midrule
\multicolumn{8}{@{}l}{\textit{Clinical DA}} \\
VRADA~\cite{purushotham2017vrada}   & \xmark & R & \cmark & \xmark & \cmark & 1 & \xmark \\
OTTEHR~\cite{ottehr2025}            & \xmark & R & \xmark & \xmark & \xmark & 3 & \xmark \\
Mutnuri et al.~\cite{mutnuri2024}   & \xmark & W & \cmark & \xmark & \cmark & 3 & \xmark \\
\midrule
\multicolumn{8}{@{}l}{\textit{General TS-DA}} \\
CoDATS~\cite{wilson2020codats}       & \xmark & R & \cmark & \xmark & \xmark & --- & \xmark \\
RAINCOAT~\cite{he2023raincoat}       & \xmark & R & \cmark & \xmark & \xmark & --- & \xmark \\
CLUDA~\cite{ozyurt2023cluda}         & \xmark & R & \cmark & \xmark & \xmark & --- & \xmark \\
\midrule
\multicolumn{8}{@{}l}{\textit{Frozen-model / Input-space}} \\
VPT~\cite{jia2022vpt}                & \cmark & I & \xmark & \xmark & \xmark & 0 & \xmark \\
Time-LLM~\cite{jin2024timellm}       & \cmark & I & \xmark & \xmark & \xmark & 0 & \xmark \\
TATO~\cite{qiu2026tato}              & \cmark & I & P      & \xmark & \cmark & 0 & \xmark \\
\midrule
\multicolumn{8}{@{}l}{\textit{Retrieval-augmented}} \\
$k$NN-MT~\cite{khandelwal2021knnmt}  & \cmark & O & \cmark & \cmark & \cmark & 0 & \xmark \\
TS-RAG~\cite{ning2025tsrag}          & \xmark & R & \cmark & \cmark & \cmark & 0 & \xmark \\
\midrule
\multicolumn{8}{@{}l}{\textit{Foundation model}} \\
ICareFM~\cite{burger2025icarefm}     & \xmark & W & \cmark & \xmark & \cmark & 4 & \xmark \\
\midrule
\textbf{Ours (Retrieval Translator)} & \cmark & \textbf{I} & \cmark & \cmark & \cmark & \textbf{5} & \cmark \\
\bottomrule
\end{tabular}
\vspace{-6pt}
\end{table}

Several observations stand out.
Representation-space (R) and weight-space (W) methods are architecturally coupled to the model's internals---they cannot be formulated without access to intermediate representations or weight tensors, violating the frozen-model constraint.
In contrast, input-space translation treats the target model as a \emph{black-box oracle} $f: \mathcal{X} \to \mathcal{Y}$.
While our efficient implementation leverages differentiable access (backpropagation through the frozen model) for gradient-based optimization, the method requires only query access to the model's input--output mapping.
In principle, this makes input-space translation applicable even to proprietary models exposed only via API, using gradient-free optimization (e.g., evolutionary strategies or Bayesian optimization)---a property no R-space or W-space method can satisfy, though we note that gradient-free training of the translator remains untested and would likely require substantially more queries.
The frozen-model methods from NLP and vision (VPT, Time-LLM) also operate in input space but use input-independent perturbations, limiting their capacity for heterogeneous clinical data.
$k$NN-MT~\cite{khandelwal2021knnmt} is closest in spirit---frozen model, retrieval,
per-timestep conditioning---but operates in output space (interpolating token distributions)
rather than input space (transforming the data itself).
Our retrieval translator is the only method that combines a fully frozen predictor,
input-space adaptation, sample-conditional neural translation, target-domain retrieval,
variable-length time-series support, evaluation across five clinical tasks, and a
theoretical gradient alignment framework.


%% ============================================================
\subsection{Head-to-Head Analysis}
\label{sec:head_to_head}

\paragraph{Against clinical DA methods.}
VRADA~\cite{purushotham2017vrada} pioneered adversarial alignment for clinical
time series but requires end-to-end training and was evaluated on a single task
(mortality).
Mutnuri et al.~\cite{mutnuri2024} address three of our five tasks (mortality,
AKI, LoS) with inductive transfer and standard adversarial DA, but fine-tune the
target model---precisely the operation our setting forbids.
On the shared AKI task, our frozen-model retrieval translator achieves
AUROC 91.14 and AUCPR 72.9, surpassing the MIMIC-native LSTM baseline (89.7 AUROC, 66.5 AUCPR) that represents
the performance ceiling when the model is trained on in-domain data---despite never modifying any predictor weights.
OTTEHR~\cite{ottehr2025} applies optimal transport to the same eICU--MIMIC
databases but performs population-level feature alignment without temporal
modeling, treating each feature independently.
Our instance-level, per-timestep retrieval and causal attention preserve the
temporal structure essential for sequential clinical tasks (AKI, sepsis, LoS).
In summary, prior clinical DA methods either sacrifice model integrity
(fine-tuning) or temporal structure (population-level alignment); we sacrifice
neither.

\paragraph{Against adapted DA baselines (DANN/CORAL/CoDATS).}
A natural question is whether standard DA methods can be adapted to our
frozen-model setting by learning an input transformation $T_\theta(\mathbf{x})$
that precedes the frozen predictor, with DANN's gradient reversal~\cite{ganin2016dann}
or CORAL's covariance alignment~\cite{sun2016coral} applied to the predictor's
internal representations.
We implement these adapted baselines (detailed in the full paper) and
find that they achieve limited improvement.
The structural argument is that clinical domain shift is \emph{heterogeneous}:
across 48 physiological features, the distributional shifts differ in magnitude,
direction, and functional form (e.g., heart rate may require a simple affine
correction while lab values exhibit nonlinear, context-dependent shifts).
A shallow input transform with a single domain discriminator or covariance
penalty lacks the capacity for such feature-specific, sample-conditional
mappings.
Our translator architectures provide structured inductive biases---additive
residuals (delta), latent bottleneck (shared latent), or target-domain cross-attention
(retrieval)---that decompose this complex mapping into learnable components.
The gradient alignment analysis (Section~\ref{sec:formulation}) further
explains the failure mode: when $\cos(\nabla_\theta \mathcal{L}_\text{task},
\nabla_\theta \mathcal{L}_\text{fid}) < 0$, as observed for sparse-label tasks,
the domain-alignment and task objectives pull the input transform in
opposing directions, a conflict that generic input transforms cannot resolve.

\paragraph{Against frozen-model methods.}
Prompt tuning~\cite{lester2021prompt} and VPT~\cite{jia2022vpt} learn a fixed
set of $\mathcal{O}(1)$ parameters (soft tokens or visual prompts) that are
shared across all inputs.
Clinical EHR time series vary from 1 to 169 timesteps, span 48--292 features with
complex missingness patterns, and exhibit patient-specific distributional
characteristics---far exceeding the capacity of input-independent perturbations.
Time-LLM~\cite{jin2024timellm} reprograms frozen LLMs for time-series forecasting,
but performs cross-\emph{modality} adaptation (language$\to$time series) for a
generic pre-trained model, not cross-\emph{domain} adaptation for a task-specific
validated clinical predictor.
TATO~\cite{qiu2026tato}, a concurrent ICLR~2026 contribution, is the closest
work in spirit: it explicitly proposes ``adapt data to model'' for frozen
time-series foundation models.
However, TATO relies on handcrafted transformation operators (slicing,
normalization, outlier correction) rather than learned neural translation,
targets forecasting rather than clinical prediction, and provides no
theoretical analysis of when data adaptation succeeds or fails.
We view TATO as validating the broader ``adapt data, not model'' paradigm;
our work extends it with (i)~learned end-to-end translation,
(ii)~retrieval-augmented cross-attention, and (iii)~gradient alignment theory.

\paragraph{Against retrieval-augmented methods.}
$k$NN-MT~\cite{khandelwal2021knnmt} is our closest architectural ancestor:
both systems pair a frozen model with per-timestep $k$-NN retrieval from a
domain-specific datastore.
The critical distinction lies in \emph{how} retrieval is used.
$k$NN-MT interpolates \emph{output token distributions}---blending the frozen
decoder's next-token prediction with the empirical distribution from retrieved
neighbors.
Our retrieval translator uses retrieved target-domain embeddings as
\emph{cross-attention context} in a Transformer decoder that transforms the
\emph{input}, enabling selective feature-level borrowing rather than
output-level interpolation.
Moreover, $k$NN-MT operates within a single language (e.g., domain-adapted
German$\to$English), while we bridge fundamentally different data distributions
with heterogeneous recording practices and device calibrations.
TS-RAG~\cite{ning2025tsrag} and RATD~\cite{liu2024ratd} retrieve time-series
segments for forecasting augmentation within the same domain; our retrieval
bridges \emph{across domains}---a qualitatively different use of the retrieval
paradigm.
RAM-EHR~\cite{xu2024ramehr} retrieves textual medical knowledge from PubMed
and knowledge graphs; we retrieve encoded \emph{time-series windows}
representing what the frozen predictor expects to see, grounding the translation
in the target domain's statistical structure rather than external knowledge.


%% ============================================================
\subsection{Results Positioning}
\label{sec:results_position}

Table~\ref{tab:results_position} situates our best results against
reference baselines from the YAIB benchmark~\cite{vandewater2024yaib}.
Full experimental details and ablations will appear in the full paper.

\begin{table}[t]
\centering
\caption{Our best translator results vs.\ reference baselines from YAIB~\cite{vandewater2024yaib}.
Classification metrics: AUROC and AUCPR $\times 100$; regression: MAE in natural units.
\textbf{Bold} indicates our result matches or surpasses the reference.
$\dagger$: MIMIC-native LSTM surpassed.
All our results use a \emph{frozen} MIMIC-IV predictor with no weight updates.}
\label{tab:results_position}
\small
\setlength{\tabcolsep}{3.2pt}
\begin{tabular}{@{}lcccccccc@{}}
\toprule
& \multicolumn{2}{c}{\textbf{Mortality}} & \multicolumn{2}{c}{\textbf{AKI}} & \multicolumn{2}{c}{\textbf{Sepsis}} & \textbf{LoS} & \textbf{KF} \\
& AUROC & AUCPR & AUROC & AUCPR & AUROC & AUCPR & MAE (h) & MAE (mg/dL) \\
\midrule
Frozen baseline  & 80.79 & 50.2 & 85.58 & 56.8 & 71.59 & 3.0 & 42.5 & 0.403 \\
eICU-native LSTM & 85.5  & 35.7 & 90.2  & 69.9 & 74.0  & 4.0 & 39.2 & 0.28 \\
MIMIC-native LSTM& 86.7  & 41.0 & 89.7  & 66.5 & 82.0  & 8.0 & 40.6 & 0.28 \\
\midrule
\textbf{Ours (best)}
  & \textbf{85.55} & \textbf{55.7} & $\mathbf{91.14}^\dagger$ & $\mathbf{72.9}^\dagger$ & \textbf{76.71} & \textbf{5.2} & \textbf{39.2} & 0.382 \\
$\Delta$ vs.\ frozen
  & $+$4.76 & $+$5.5 & $+$5.56 & $+$16.1 & $+$5.12 & $+$2.2 & $-$3.3 & $-$0.021 \\
\bottomrule
\end{tabular}
\vspace{-4pt}
\end{table}

\paragraph{Where we exceed expectations.}
On AKI, our frozen-model translator achieves AUROC 91.14, surpassing not only
the eICU-native LSTM (90.2) but also the \emph{MIMIC-native LSTM} (89.7)---the
model trained on in-domain data with full weight optimization.
The AUCPR results are even more striking: $72.9$ vs.\ $66.5$ for the MIMIC-native LSTM and $69.9$ for the eICU-native LSTM, a $+16.1$ point improvement over the frozen baseline.
AUCPR is particularly informative for clinical tasks with class imbalance because it measures the precision--recall tradeoff directly: a higher AUCPR means the model identifies more true positives without generating excessive false alarms.
The fact that the AUCPR gain ($+16.1$) is nearly $3\times$ the AUROC gain ($+5.6$) indicates that translation disproportionately improves the model's ability to distinguish positive cases at clinically relevant operating points.
On mortality, translation pushes AUCPR to $55.7$---far above both the eICU-native LSTM ($35.7$) and the MIMIC-native LSTM ($41.0$).
Notably, the frozen baseline already exceeds both native models on mortality AUCPR ($50.2$), and the translator extends this advantage further.
This demonstrates that input-space translation can recover performance
\emph{beyond} what the target model achieves on its own training distribution,
suggesting that the translated source data may introduce beneficial distributional
diversity that complements the frozen predictor's training distribution.
All three classification tasks surpass the eICU-native LSTM on both AUROC and AUCPR, confirming that
the translator closes the domain gap rather than merely compensating for it.
On LoS (regression), the translator achieves MAE of 39.2~hours, matching the
eICU-native LSTM exactly and improving upon the MIMIC-native LSTM (40.6~hours).
Cross-server reproducibility is strong: across 3 seeds on 2 physically
distinct GPU servers, AKI AUROC varies by only $\pm$0.0005, confirming that
the gains are robust and not artifacts of hardware-specific numerical behavior.

\paragraph{Where we acknowledge gaps.}
We identify three areas where performance does not yet match the strongest
reference points, and in each case, we characterize the underlying cause.
\emph{Sepsis}: the retrieval translator achieves $+5.12$ AUROC and $+2.2$ AUCPR, surpassing the eICU-native LSTM on both metrics (76.71 vs.\ 74.0 AUROC; 5.2 vs.\ 4.0 AUCPR---a 30\% relative AUCPR improvement), but the gains are smaller in magnitude than on denser tasks.
With only 1.13\% positive labels per timestep, the task gradient is extremely sparse: our gradient alignment analysis ($\alpha = -0.21$ for sepsis vs.\ $+0.84$ for mortality) correctly predicts the relative difficulty and identifies label sparsity as the primary bottleneck for further improvement.
\emph{Kidney Function}: the KF task involves 292 features including 192
generated cumulative statistics (\texttt{\_min\_hist}, \texttt{\_max\_hist},
\texttt{\_mean\_hist}, \texttt{\_count}) that must maintain internal
consistency after translation---a substantially harder translation target.
Our MAE of 0.382 improves over the frozen baseline (0.403) but a gap to the
eICU-native reference (0.28) remains.
\emph{Single domain pair}: our current evaluation covers eICU$\to$MIMIC-IV;
extension to HiRID$\to$MIMIC-IV is ongoing and will test generalization of
the approach across a third institutional distribution.


%% ============================================================
\subsection{Contributions}
\label{sec:contributions}

Our work makes five concrete contributions:

\begin{enumerate}[leftmargin=*,itemsep=2pt]

\item \textbf{Frozen-model preservation as a first-class constraint.}
We formalize the setting where the \emph{entire} target predictor is frozen---not
just the classifier head (SHOT~\cite{liang2020shot}) or a subset of layers---and
show that meaningful domain adaptation is achievable purely through input-space
translation.
This is directly motivated by the FDA's Software as a Medical Device (SaMD)
framework~\cite{fda2021samd}, which provides a simpler regulatory pathway for
``locked'' algorithms whose weights are not modified post-deployment.

\item \textbf{Input-space modularity.}
The translator is an external, independently trainable module that can be
swapped, updated, or audited without touching the validated predictor.
This decoupling enables a deployment model where the predictor and translator
have separate validation lifecycles---a practical requirement in regulated
clinical environments.

\item \textbf{Instance-level retrieval-augmented cross-attention for clinical
time-series DA.}
To the best of our knowledge, we are the first to combine per-timestep $k$-NN
retrieval from a target-domain memory bank with cross-attention decoding for
domain adaptation of clinical time series.
This provides each translated timestep with contextual grounding in what the
frozen model expects to see, enabling the fine-grained, feature-specific
transformations that generic input transforms cannot achieve.

\item \textbf{Universal paradigm across five clinical tasks.}
We evaluate on the broadest task set in the clinical DA literature:
three classification tasks (Mortality24, AKI, Sepsis) spanning per-stay and
per-timestep label structures, and two regression tasks (Length of Stay, Kidney
Function) with continuous targets.
The retrieval translator achieves the best or tied-best results on all five,
with improvements on both AUROC and AUCPR for all classification tasks---including
surpassing the MIMIC-native LSTM on both metrics for AKI.
This establishes retrieval translation as a universal paradigm with task-adapted capacity
(specifically, the number of cross-attention layers: 3 for dense-label tasks
like AKI, 2 for sparse-label tasks like sepsis).

\item \textbf{Gradient alignment as a predictive diagnostic.}
We introduce the cosine similarity
$\cos(\nabla_\theta \mathcal{L}_\text{task}, \nabla_\theta \mathcal{L}_\text{fid})$
as a diagnostic for frozen-model translation, connecting to the gradient conflict
literature in multi-task learning~\cite{yu2020pcgrad,liu2021cagrad} and
UDA~\cite{du2024gh,phan2024pga}.
This metric correctly predicts method--task suitability a priori
($+0.84$ for mortality, where delta translation succeeds; $-0.21$ for sepsis,
where it fails and retrieval is required) and reveals that the retrieval
architecture resolves the conflict \emph{structurally} through the cross-attention
information pathway, rather than requiring algorithmic gradient surgery.

\end{enumerate}


%% ============================================================
\subsection{Scope, Applicability, and Limitations}
\label{sec:limitations}

\paragraph{When to use input-space translation.}
Our approach is most applicable when three conditions hold:
(i)~a validated target-domain predictor exists and must remain frozen
(regulatory, institutional, or practical constraints);
(ii)~labeled data from both source and target domains is available for
translator training; and
(iii)~the target task provides sufficient label signal for the task gradient to guide translation. In our experiments, translation gains scale with label density: AKI (11.95\%) and mortality (5.5\%) yield the largest improvements, while sepsis (1.13\%) shows smaller but still meaningful gains ($+5.12$ AUROC), consistent with the gradient alignment analysis.
Under these conditions, the retrieval translator consistently closes the
domain gap across diverse clinical prediction tasks.

\paragraph{When alternatives are preferable.}
If the target model can be retrained, standard domain adaptation or foundation
model fine-tuning can optimize the predictor's own parameters---a strictly more
expressive search space---and may outperform input-space translation, though
this is not guaranteed: our translator surpasses the in-domain MIMIC-native
LSTM on AKI, suggesting that the inductive bias of retrieval-augmented
translation can sometimes compensate for the constrained parameter space.
If no target-domain data is available, source-free methods
(e.g., MAPU~\cite{ragab2023mapu}) are one option; alternatively,
our delta translator can operate without a target-domain memory bank,
using only the frozen model's task gradient and fidelity regularisation
to learn input corrections---achieving up to $+2.9$ AUROC on AKI
(vs.\ $+5.6$ with retrieval) when gradient alignment is positive ($\alpha > 0$).
If real-time inference latency is critical, the memory bank $k$-NN lookup
introduces approximately $2$--$3\times$ overhead compared to the frozen model
alone, which may be unacceptable in time-sensitive clinical alerting systems.

\paragraph{Current limitations.}
We identify five limitations that define the scope of the present work and
inform future directions:
\begin{enumerate}[leftmargin=*,itemsep=1pt]
\item \textit{Single domain pair.}
Our evaluation covers eICU$\to$MIMIC-IV.
Extension to additional source domains (HiRID$\to$MIMIC-IV) is under way
and will test whether the approach generalizes across institutional
distributions with different data collection protocols.

\item \textit{Label density dependence.}
Translation quality correlates with label density:
AKI (11.95\% positive rate) shows the largest gains ($\Delta$AUROC $+$5.56, $\Delta$AUCPR $+$16.1),
while sepsis (1.13\%) shows smaller but still meaningful improvements ($+5.12$ AUROC, $+2.2$ AUCPR).
The gradient alignment diagnostic quantifies this relationship, but does
not resolve the underlying information bottleneck for extremely sparse tasks.

\item \textit{Calibration degradation.}
Translation improves discrimination (AUROC, AUCPR) but can degrade probability
calibration (ECE), requiring post-hoc temperature scaling---an additional
validation step in clinical deployment.

\item \textit{Inference overhead.}
The retrieval translator requires maintaining a GPU-resident memory bank and
performing $k$-NN search at each timestep, adding computational and memory
costs relative to direct model inference.

\item \textit{Architecture specificity.}
Our frozen baselines are LSTMs from the YAIB
framework~\cite{vandewater2024yaib}.
Since input-space translation treats the predictor as a black-box function
$f: \mathcal{X} \to \mathcal{Y}$, the paradigm is architecture-agnostic in
principle; empirical validation with other target architectures (Transformers,
GRUs, temporal convolutional networks) remains future work.
\end{enumerate}

\noindent
We view these as open problems rather than fundamental obstacles.
The frozen-model, input-space translation paradigm opens a new direction in
clinical DA research, and the gradient alignment framework provides guidance
for extending it to new domains, tasks, and architectures.
